\documentclass[12pt, letterpaper]{article}

\usepackage[utf8]{inputenc}
\usepackage[T1]{fontenc}
\usepackage[spanish]{babel}
\usepackage{geometry}
\usepackage{amsmath, amssymb, amsthm}
\usepackage{booktabs}
\usepackage{array}
\usepackage{microtype}
\usepackage{setspace}
\usepackage{parskip}
\usepackage{titlesec}
\usepackage{fancyhdr}
\usepackage{enumitem}
\usepackage{bm}
\usepackage{mathtools}
\usepackage{tcolorbox}
\usepackage{hyperref}

\tcbuselibrary{skins, breakable}

\geometry{top=2.8cm, bottom=3.0cm, left=3.2cm, right=3.2cm, headheight=14pt}
\setstretch{1.20}

\pagestyle{fancy}
\fancyhf{}
\fancyhead[C]{\small\textsc{Econometria · Gujarati --- Ejercicio 12.2}}
\fancyfoot[C]{\small\thepage}
\renewcommand{\headrulewidth}{0.4pt}
\renewcommand{\footrulewidth}{0pt}

\titleformat{\section}
  {\large\bfseries\scshape}{\thesection.}{0.6em}{}
  [\vspace{-4pt}\rule{\linewidth}{0.4pt}]
\titleformat{\subsection}{\normalsize\bfseries}{\thesubsection.}{0.5em}{}
\titleformat{\subsubsection}{\normalsize\itshape}{}{0em}{}

\newtcolorbox{clave}{
  colback=white, colframe=black,
  boxrule=0.5pt, arc=3pt,
  left=8pt, right=8pt, top=6pt, bottom=6pt,
  breakable
}

\newcommand{\E}{\mathbb{E}}
\newcommand{\Var}{\operatorname{Var}}

\begin{document}

\begin{center}
  {\LARGE\bfseries Prueba de Durbin-Watson:}\\[6pt]
  {\LARGE\bfseries Analisis de Autocorrelacion ($n=50$, $k'=4$)}\\[12pt]
  {\large\itshape Aplicacion de las reglas de decision del estadistico $d$}\\
  {\large\itshape para distintos valores muestrales}\\[14pt]
  \rule{0.55\linewidth}{0.6pt}\\[8pt]
  {\normalsize Ejercicio~12.2 --- \textit{Econometria}, Gujarati \& Porter, 5.a ed.\ por Emanuel Quintana Silva}\\[4pt]
  {\small\textsc{Valores criticos, zonas de decision y conclusiones de inferencia}}
\end{center}

\vspace{10pt}

% ─────────────────────────────────────────────────────────────────────────────
\section{Valores criticos de la prueba}
% ─────────────────────────────────────────────────────────────────────────────

Para $n = 50$ observaciones y $k' = 4$ variables explicativas (excluyendo la
constante), los valores criticos de la tabla de Durbin-Watson al nivel de
significancia del $5\%$ son:
\begin{equation}
  d_L = 1.378, \qquad d_U = 1.721
  \label{eq:criticos}
\end{equation}

La prueba cubre cinco regiones de decision:

\begin{center}
\renewcommand{\arraystretch}{1.3}
\small
\begin{tabular}{llp{6cm}}
\toprule
\textbf{Region} & \textbf{Criterio} & \textbf{Decision} \\
\midrule
I   & $d < d_L = 1.378$
    & Rechazar $H_0$: autocorrelacion positiva \\
II  & $d_L \le d \le d_U$, \; $1.378 \le d \le 1.721$
    & Zona de indecision \\
III & $d_U < d < 4 - d_U$, \; $1.721 < d < 2.279$
    & No rechazar $H_0$: sin autocorrelacion \\
IV  & $4-d_U \le d \le 4-d_L$, \; $2.279 \le d \le 2.622$
    & Zona de indecision (autocorrelacion negativa) \\
V   & $d > 4 - d_L = 2.622$
    & Rechazar $H_0$: autocorrelacion negativa \\
\bottomrule
\end{tabular}
\end{center}

\medskip
\noindent Para contrastar autocorrelacion \emph{negativa} se utiliza el estadistico
transformado $d^* = 4 - d$, que se compara con los mismos valores criticos $d_L$
y $d_U$.

% ─────────────────────────────────────────────────────────────────────────────
\section{Analisis de cada valor del estadistico $d$}
% ─────────────────────────────────────────────────────────────────────────────

\subsection{Inciso (a): $d = 1.05$}

\textbf{Comparacion:}
\begin{equation}
  d = 1.05 < d_L = 1.378
\end{equation}

El valor observado se encuentra en la \textbf{Region~I}, por debajo del limite
inferior. Se rechaza la hipotesis nula de no autocorrelacion.

\textbf{Conclusion:} Existe evidencia estadistica de \textbf{autocorrelacion
positiva} de primer orden al nivel de significancia del $5\%$. Los residuos MCO
muestran un patron de correlacion serial positiva, lo que implica que errores
consecutivos tienden a tener el mismo signo.

\begin{clave}
  $d = 1.05 < d_L = 1.378$ $\Rightarrow$ \textbf{Rechazar $H_0$}.
  Autocorrelacion positiva significativa.
\end{clave}

\subsection{Inciso (b): $d = 1.40$}

\textbf{Comparacion:}
\begin{equation}
  d_L = 1.378 \;\le\; d = 1.40 \;\le\; d_U = 1.721
\end{equation}

El valor observado cae exactamente en la \textbf{Region~II} (zona de indecision).
La prueba de Durbin-Watson no permite concluir si existe o no autocorrelacion
positiva con la informacion disponible.

\textbf{Conclusion:} La prueba es \textbf{inconcluyente}. En la practica
econometrica se recomienda complementar el analisis con la prueba de
Breusch-Godfrey (prueba LM), que no presenta zonas de indecision y es valida
incluso con modelos dinamicos.

\begin{clave}
  $d_L \le 1.40 \le d_U$ $\Rightarrow$ \textbf{Zona de indecision}.
  Usar prueba de Breusch-Godfrey como alternativa.
\end{clave}

\subsection{Inciso (c): $d = 2.50$}

Para contrastar autocorrelacion negativa se calcula:
\begin{equation}
  d^* = 4 - d = 4 - 2.50 = 1.50
\end{equation}

\textbf{Comparacion respecto a autocorrelacion negativa:}
\begin{equation}
  d_L = 1.378 \;\le\; d^* = 1.50 \;\le\; d_U = 1.721
  \quad \Rightarrow \quad \text{zona de indecision}
\end{equation}

\textbf{Comparacion respecto a autocorrelacion positiva:}
\begin{equation}
  d = 2.50 > d_U = 1.721
  \quad \Rightarrow \quad \text{no hay autocorrelacion positiva}
\end{equation}

Equivalentemente, los limites para la region de indecision de autocorrelacion
negativa son $4 - d_U = 2.279$ y $4 - d_L = 2.622$, y se verifica:
\begin{equation}
  2.279 \;\le\; d = 2.50 \;\le\; 2.622
\end{equation}

\textbf{Conclusion:} La prueba es \textbf{inconcluyente} respecto a la
autocorrelacion negativa (Region~IV). No hay evidencia de autocorrelacion
positiva.

\begin{clave}
  $d = 2.50$: inconcluyente para autocorrelacion negativa;
  sin autocorrelacion positiva.
\end{clave}

\subsection{Inciso (d): $d = 3.97$}

Se calcula el estadistico transformado:
\begin{equation}
  d^* = 4 - d = 4 - 3.97 = 0.03
\end{equation}

\textbf{Comparacion:}
\begin{equation}
  d^* = 0.03 \;<\; d_L = 1.378
\end{equation}

El estadistico transformado cae en la Region~I del criterio para
autocorrelacion negativa, equivalente a la Region~V del estadistico original
($d = 3.97 > 4 - d_L = 2.622$). Se rechaza la hipotesis nula.

\textbf{Conclusion:} Existe evidencia estadistica muy fuerte de
\textbf{autocorrelacion negativa} de primer orden al nivel del $5\%$. Los
errores consecutivos tienden a alternar de signo, lo que puede indicar una
sobrecoleccion de errores o la presencia de un proceso de medias moviles en los
residuos.

\begin{clave}
  $d^* = 4 - 3.97 = 0.03 < d_L = 1.378$ $\Rightarrow$ \textbf{Rechazar $H_0$}.
  Autocorrelacion negativa significativa.
\end{clave}

% ─────────────────────────────────────────────────────────────────────────────
\section{Sintesis de resultados}
% ─────────────────────────────────────────────────────────────────────────────

\begin{center}
\renewcommand{\arraystretch}{1.4}
\small
\begin{tabular}{cccl}
\toprule
\textbf{Inciso} & \textbf{$d$ observado} & \textbf{Region} & \textbf{Conclusion} \\
\midrule
(a) & $1.05$ & I              & Autocorrelacion positiva significativa \\
(b) & $1.40$ & II             & Inconcluyente (zona de indecision) \\
(c) & $2.50$ & IV             & Inconcluyente para autocorrelacion negativa \\
(d) & $3.97$ & V              & Autocorrelacion negativa significativa \\
\bottomrule
\end{tabular}
\end{center}

\vspace{14pt}
\noindent\rule{\linewidth}{0.4pt}

\noindent{\small\textit{Nota.}
La zona de indecision es la limitacion mas conocida de la prueba de
Durbin-Watson. Cuando el estadistico cae en ella, la alternativa mas utilizada
en econometria aplicada es la \textbf{prueba de Breusch-Godfrey} (o prueba LM
de multiplicadores de Lagrange): no tiene zonas de indecision, es valida con
regresores estocasticos y permite contrastar autocorrelacion de orden superior
al primero. Para series temporales con posible raiz unitaria en el error,
conviene ademas aplicar la prueba de \textbf{Berenblut-Webb} ($g$-test), que
contrasta especificamente $H_0: \rho = 1$.}

\end{document}
